% Part: first-order-logic
% Chapter: introduction
% Section: semantic-notions

\documentclass[../../../include/open-logic-section]{subfiles}

\begin{document}

\olfileid{fol}{int}{sem}

\olsection{معنوی تصورات}

اسیں اُتّے ذکر کیتا سی کہ جدوں اسیں ایہ ویکھدے آں کہ
$\Sat{M}{!A}[s]$ قائم اے یا نہیں، تاں سہولت لئی
$s$ نوں سارے !!{variable}s نوں قدراں دین دیندے آں؛
پر ورتیاں صرف اوہ قدراں جاندیاں نیں جیہڑیاں اوہ
~$!A$ وچ آؤن والے !!{variable}s نوں دیندا اے۔
اصل وچ~$!A$ دے صرف \emph{آزاد} متغیراں دیاں
قدراں اہم نیں۔ اسیں احتیاط نال کم کردے آں،
ایس لئی ایہ گل ثابت وی کراں گے۔ چونکہ
!!{sentence}s وچ کوئی آزاد متغیر نہیں ہوندا،
!!a{structure} وچ انہاں دے پورا اترن یا نہ
اترن لئی $s$ دا کوئی اثر نہیں۔ ایس لئی
جدوں~$!A$ اک !!a{sentence} ہووے، تاں
$\Sat{M}{!A}$ دی تعریف ایہ دے سکدے آں کہ
”$\Sat{M}{!A}[s]$ سارے~$s$ لئی“؛
تے ایہ اُدوں تے صرف اُدوں سچ اے جدوں
$\Sat{M}{!A}[s]$ گھٹ توں گھٹ اک~$s$ لئی ہووے۔
!!{formula}s دے پورا اترن دی قابلِ عمل تعریف
لئی سانوں !!{variable} قدر-بندیاں متعارف
کرانیاں پئیاں، پر !!{sentence}s لئی پورا
اترنا !!{variable} قدر-بندیاں توں آزاد اے۔

جدوں ”$\Sat{M}{!A}$“ دی تعریف ساڈے کول آ
جاندی اے، تاں پہلے درجے دی منطق دے
!!{sentence}s بارے ”حالت“ تے ”وچ سچا ہون“
دا مطلب وی واضح ہو جاندا اے۔ !!{sentence}s
لئی $\Sat{M}{!A}$ دی تعریف دی بنیاد اُتے
فیر اسیں منطقی اعتبار، منطقی نتیجہ تے پورا
ہون جوگتا دے بنیادی معنوی تصورات تعریف
کر سکدے آں۔ اک جملہ معتبر اے،
$\Entails !A$، جے ہر !!{structure}
اوہنوں پورا کردی اے۔ !!{sentence}s دے کسے سیٹ
توں اوہ منطقی طور تے نکلدا اے،
$\Gamma \Entails !A$، جے ہر اوہ
!!{structure} جیہڑی~$\Gamma$ دے سارے
!!{sentence}s نوں پورا کردی اے،
~$!A$ نوں وی پورا کردی ہووے۔ تے
!!{sentence}s دا اک سیٹ پورا ہون جوگا
اے جے کوئی !!{structure} اوہدے سارے
!!{sentence}s نوں اکّو ویلے پورا کردی ہووے۔

چونکہ !!{formula}s دی تعریف استقرائی اے،
تے پورا اترن دی تعریف وی !!{formula}s دی
بناوٹ اُتے استقرا نال دتی جاندی اے،
اسیں استقرا ورت کے اپنی معنویات دیاں
خاصیتاں ثابت کر سکدے آں تے تعریف شدہ
معنوی تصورات دا آپسی تعلق وی واضح
کر سکدے آں۔ اسیں ایہناں وچوں کجھ
خاصیتاں اکٹھیاں کر کے ثابت کراں گے:
کجھ ایس لئی کہ اوہ آپ دلچسپ نیں،
پر زیادہ ایس لئی کہ جدوں اسیں معنویات
تے !!{derivation} نظاماں دے رشتے دی
جانچ کراں گے، تاں ایہ کم آؤن گیاں۔
ایس کم لئی سانوں کجھ اوہ نحوی تصورات
تے عمل وی ٹھیک ٹھیک، یعنی استقرائی
طور تے، تعریف کرنے پین گے جنہاں دا
ہن تک ذکر نہیں کیتا۔

\end{document}
